Augmented matrices can be used to find solutions to systems of linear equations by \emph{reducing} them using the elementary row operations.

\begin{definition}
The first nonzero entry in a nonzero row is called the \emph{leading entry.} A matrix is \emph{reduced} if
\begin{enumerate}
    \item All zero rows are at the bottom.
    \item For each nonzero row, the leading entry is one and all other entries in the column are zero.
    \item The entry in each row is to the right of the leading entry in any row above it.
\end{enumerate}
\end{definition}

\begin{example}
    The following are reduced matrices.
    
    \[\begin{bmatrix}
        1 & 0 & 3 \\
        0 & 1 & 1
    \end{bmatrix}\]

    Note that if we take this matrix as an augmented matrix that represents a system of equations, we would have
    \[\begin{cases}
        x = 3 \\
        y = 1
    \end{cases}\]
    and the solution to this system is immediately apparent: \(\left(3,1\right).\)
\end{example}
\begin{example}

    \[\begin{bmatrix}
        \colorbox{cyan}{1} & 0 & 2 & 0 \\
        0 & \colorbox{cyan}{1} & 6 & 0 \\
        0 & 0 & 0 & \colorbox{cyan}{1}
    \end{bmatrix}\]

    This is a reduced matrix with the leading entries highlighted. There is no leading entry in the third column. Thus we do not care about the nonzero entries in that column.

    As a system of equations,
    \[\begin{cases}
        x + 2z = 0 \\
        y + 6z = 0 \\
        0=1
    \end{cases}.\]
    Given that \(0=1\) has no solutions, this system has \emph{no solutions.} When reducing a matrix to solve a system of a equations, if we end up with a leading entry in the last column, there is no need to continue reducing, we can immediately conclude that there is no solution.
\end{example}

We reduce matrices to solve systems of equations by applying the elementary row operations. This method is analogous to the method of solving a system by elimination.

When performing the following row operations on an augmented matrix, the set of solutions is preserved.

THE ROW OPERATIONS
\begin{enumerate}
    \item Exchanging rows.
    \item Multiplying a row by a scalar.
    \item Replacing a row with the sum of that row with another.
\end{enumerate}

\begin{example}
    Find all solutions to the system of linear equations.
   \[ \begin{cases}
        2x + y - z - 1 = 0\\
        x + 2y +2z -3 = 0 \\
        y + 5z - 2 =0
    \end{cases}\]
    \begin{solution}
        We first set up the augmented matrix.
        \[\begin{bmatrix}
            2 & 1 & -1 & \vert & 1 \\
            1 & 2 & 2 & \vert & 3 \\
            0 & 1 & 5 & \vert & 2
        \end{bmatrix}\]

        We will let \(R_1, R_2,\) and \(R_3\) be the first, second, and third rows respectively.
        We first go down the columns from left to right making the non leading entries below each leading entry zero. After this we go up the columns from right to left making the necessary entries zero.

        \begin{gather*}\overset{R_1/2}{\longrightarrow} \begin{bmatrix}
            1 & 1/2 & -1/2 & \vert & 1/2 \\
            1 & 2 & 2 & \vert & 3 \\
            0 & 1 & 5 & \vert & 2
        \end{bmatrix} \\
        \overset{R_2-R_1}{\longrightarrow} \begin{bmatrix}
            1 & 1/2 & -1/2 & \vert & 1/2 \\
            0 & 3/2 & 5/2 & \vert & 5/2 \\
            0 & 1 & 5 & \vert & 2
        \end{bmatrix}\\
        \overset{2R_2/3}{\longrightarrow} \begin{bmatrix}
            1 & 1/2 & -1/2 & \vert & 1/2 \\
            0 & 1 & 5/3 & \vert & 5/3 \\
            0 & 1 & 5 & \vert & 2
        \end{bmatrix}\\
        \overset{R_3-R_2}{\longrightarrow} \begin{bmatrix}
            1 & 1/2 & -1/2 & \vert & 1/2 \\
            0 & 1 & 5/3 & \vert & 5/3 \\
            0 & 0 & 10/3 & \vert & 1/3
        \end{bmatrix}\\
        \overset{3R_3/10}{\longrightarrow} \begin{bmatrix}
            1 & 1/2 & -1/2 & \vert & 1/2 \\
            0 & 1 & 5/3 & \vert & 5/3 \\
            0 & 0 & 1 & \vert & 1/10
        \end{bmatrix}
        \end{gather*}
        We now have a \(1\) as a leading entry for each row, as well as \(0\)s in all of the entries below each leading entry. We now want to clear out all of the entries above the leading \(1\)s.
        \begin{gather*}
        \overset{R_2 - 5R_3/3}{\longrightarrow} \begin{bmatrix}
            1 & 1/2 & -1/2 & \vert & 1/2 \\
            0 & 1 & 0 & \vert & 3/2 \\
            0 & 0 & 1 & \vert & 1/10
        \end{bmatrix}\\
        \overset{R_1 + R_3/2}{\longrightarrow} \begin{bmatrix}
            1 & 1/2 & 0 & \vert & 11/20 \\
            0 & 1 & 0 & \vert & 3/2 \\
            0 & 0 & 1 & \vert & 1/10
        \end{bmatrix}\\
        \overset{R_1 + R_3/2}{\longrightarrow} \begin{bmatrix}
            1 & 0 & 0 & \vert & -1/5 \\
            0 & 1 & 0 & \vert & 3/2 \\
            0 & 0 & 1 & \vert & 1/10
        \end{bmatrix}
        \end{gather*}
        Now that we have a reduced matrix, we can now determine the solution to the system of equations by converting this matrix back into a system of linear equations.
        \[\begin{cases}
            x = -1/5 \\
            y = 3/2 \\
            z = 1/10
        \end{cases}\]
    \end{solution}
\end{example}

\begin{example}
    Find all solutions to the system of linear equations.
    \[\begin{cases}
        2x_1 + 3x_2 = 5\\
        x_1 + 2x_2 = 1 \\
        x_1+5x_2=3
    \end{cases}\]
    \begin{solution}

        This system as an augmented matrix.

        \[\begin{bmatrix}
            2 & 3 & \vert & 5 \\
            1 & 2 & \vert & 1 \\
            1 & 5 & \vert & 3
        \end{bmatrix}\]
        We will let \(R_1, R_2,\) and \(R_3\) be the first, second, and third rows respectively.

        \begin{gather*}\overset{R_1/2}{\longrightarrow} \begin{bmatrix}
            1 & 3/2 & \vert & 5/2 \\
            1 & 2 & \vert & 1 \\
            1 & 5 & \vert & 3
        \end{bmatrix} \\
        \overset{R_2-R_1 \ \& \ R_3-R_1}{\longrightarrow} \begin{bmatrix}
            1 & 3/2 & \vert & 5/2 \\
            0 & 1/2 &  \vert & -3/2 \\
            0 & 7/2 &  \vert & 1/2
        \end{bmatrix}\\
        \overset{2R_2 \ \& \ 2R_3/7-2R_1}{\longrightarrow} \begin{bmatrix}
            1 & 3/2 & \vert & 5/2 \\
            0 & 1 &  \vert & -3 \\
            0 & 0 &  \vert & 22/7
        \end{bmatrix}
        \end{gather*}
        We stop here because we have a leading entry in the last row. Converting this to an equation we get \(0x_1+0x_2 = 22/7.\) There are \emph{no} solutions to this equation and hence there are no solutions to the \emph{entire} system of equations.
    \end{solution}
\end{example}